\documentclass{beamer}

\usepackage{graphicx}
\usepackage{xcolor}
\usepackage{tikz}
\usepackage{amsmath,amssymb}

% ---------------- Colors ----------------
\definecolor{purpleblue}{RGB}{80, 80, 180}
\definecolor{lightgray}{RGB}{230, 230, 230}
\definecolor{novelPurple}{RGB}{98, 54, 255}
\definecolor{darkgray}{RGB}{50, 50, 50}

% ---------------- Theme ----------------
\usetheme{Madrid}
\usecolortheme{dolphin}

\setbeamerfont{title}{size=\Large,series=\bfseries}
\setbeamerfont{frametitle}{size=\LARGE,series=\bfseries}

\setbeamercolor{frametitle}{fg=white,bg=novelPurple}
\setbeamercolor{title}{fg=white,bg=purpleblue}

% ---------------- Title ----------------
\title[SAT Hard Question]{\textbf{SAT Hard Question 6}}
\author{\textbf{Jack Zeng}}
\institute{\textcolor{white}{\textbf{Novel Prep}}}
\date{\today}

\begin{document}

% ---------------- Title Page ----------------
\begin{frame}
    \titlepage
    \vfill
    \begin{flushright}
        \textcolor{purpleblue}{\Huge \textbf{Novel Prep}}
    \end{flushright}
\end{frame}

% =========================================================
% Question 1
% =========================================================

\begin{frame}{Question 1}

Given the expression

\[
\sqrt[5]{70n\left(\sqrt[5]{70n}\right)^2}
\]

for what value of $x$ is the given expression equivalent to $(70n)^{30x}$, where $n>1$?

\end{frame}

\begin{frame}{Answer to Question 1}

\textbf{Correct Answer: $\frac{3}{350}$}

\end{frame}

% =========================================================
% Question 2
% =========================================================

\begin{frame}{Question 2}

The expression $4x^2+bx-45$, where $b$ is a constant, can be rewritten as

\[
(hx+k)(x+j)
\]

where $h$, $k$, and $j$ are integer constants.

Which of the following must be an integer?

\begin{itemize}
    \item[A.] $\frac{b}{h}$
    \item[B.] $\frac{b}{k}$
    \item[C.] $\frac{45}{h}$
    \item[D.] $\frac{45}{k}$
\end{itemize}

\end{frame}

\begin{frame}{Answer to Question 2}

\textbf{Correct Answer: D}

\end{frame}

% =========================================================
% Question 3
% =========================================================

% =========================================================
% Question 3
% =========================================================

\begin{frame}{Question 3}

The expression

\[
6\sqrt[5]{3x^{45}}\cdot \sqrt[8]{8x}
\]

is equivalent to $ax^b$, where $a$ and $b$ are positive constants and $x>1$.

What is the value of $b$?

\end{frame}

\begin{frame}{Answer to Question 3}

\textbf{Correct Answer: $\frac{73}{8}$}

\end{frame}

% =========================================================
% Question 4
% =========================================================

\begin{frame}{Question 4}

In the given equation, $c$ is a positive constant.

Solve for $x$:

\[
\frac{x^2}{\sqrt{x^2-c^2}}
=
\frac{c^2}{\sqrt{x^2-c^2}}+39
\]

Which of the following is one of the solutions to the given equation?

\begin{itemize}
    \item[A.] $-c$
    \item[B.] $-c^2-39^2$
    \item[C.] $-\sqrt{39^2-c^2}$
    \item[D.] $\sqrt{c^2+39^2}$
\end{itemize}

\end{frame}

\begin{frame}{Answer to Question 4}

\textbf{Correct Answer: D}

\end{frame}
% =========================================================
% Question 5
% =========================================================

\begin{frame}{Question 5}

The quadratic function is given as

\[
f(x)=ax^2+4x+c
\]

where $a$ and $c$ are constants.

The graph of $y=f(x)$ is a parabola that opens upward and has a vertex at $(h,k)$, where $h$ and $k$ are constants.

If $k<0$ and $f(-9)=f(3)$, which of the following must be true?

\begin{itemize}
    \item[I.] $c<0$
    \item[II.] $a\ge 1$
\end{itemize}

\begin{itemize}
    \item[A.] I only
    \item[B.] II only
    \item[C.] I and II
    \item[D.] Neither I nor II
\end{itemize}

\end{frame}

\begin{frame}{Answer to Question 5}

\textbf{Correct Answer: D}

\end{frame}

% =========================================================
% Question 6
% =========================================================

\begin{frame}{Question 6}

The functions $f$ and $g$ are defined by the given equations, where $x\ge0$:

\[
f(x)=33(0.4)^{x+3}
\]

\[
g(x)=33(0.16)(0.4)^{x-2}
\]

Which of the following equations displays, as a constant or coefficient, the maximum value of the function it defines, where $x\ge0$?

\begin{itemize}
    \item[A.] I only
    \item[B.] II only
    \item[C.] I and II
    \item[D.] Neither I nor II
\end{itemize}

\end{frame}

\begin{frame}{Answer to Question 6}

\textbf{Correct Answer: B}

\end{frame}

% =========================================================
% Question 7
% =========================================================

\begin{frame}{Question 7}

Each of the following frequency tables represents a data set.

Which of these frequency tables represents the data set with the smallest standard deviation?

\vspace{0.3cm}

\begin{tabular}{c|c|c|c|c}
Value & Frequency (A) & Frequency (B) & Frequency (C) & Frequency (D)\\
\hline
5 & 0 & 11 & 15 & 18\\
6 & 10 & 17 & 15 & 14\\
7 & 55 & 19 & 15 & 11\\
8 & 10 & 17 & 15 & 14\\
9 & 0 & 11 & 15 & 18
\end{tabular}

\vspace{0.3cm}

\begin{itemize}
    \item[A.] A
    \item[B.] B
    \item[C.] C
    \item[D.] D
\end{itemize}

\end{frame}

\begin{frame}{Answer to Question 7}

\textbf{Correct Answer: A}

\end{frame}

% =========================================================
% Question 8
% =========================================================

\begin{frame}{Question 8}

In $\triangle ABC$, $AB=4$ and $BC=5$.

If $\angle ABC<90^\circ$, which of the following could be the area of $\triangle ABC$?

\begin{itemize}
    \item[A.] 5
    \item[B.] 10
    \item[C.] 15
    \item[D.] 20
\end{itemize}

\end{frame}

\begin{frame}{Answer to Question 8}

\textbf{Correct Answer: A}

\end{frame}

% =========================================================
% Question 9
% =========================================================

\begin{frame}{Question 9}

The table summarizes the distribution of heights, in inches (“in”), for 300 participants in a study.

If a participant that is at most 71 inches tall is chosen at random, what is the probability that the participant was selected from Group 2?

\vspace{0.3cm}

\begin{tabular}{c|c|c|c|c}
 & 48--59 in. & 60--71 in. & 72--83 in. & Total\\
\hline
Group 1 & 20 & 60 & 20 & 100\\
Group 2 & 10 & 75 & 15 & 100\\
Group 3 & 30 & 45 & 25 & 100\\
\hline
Total & 60 & 180 & 60 & 300
\end{tabular}

\vspace{0.3cm}

\begin{itemize}
    \item[A.] $\frac14$
    \item[B.] $\frac{17}{48}$
    \item[C.] $\frac{17}{60}$
    \item[D.] $\frac{17}{20}$
\end{itemize}

\end{frame}

\begin{frame}{Answer to Question 9}

\textbf{Correct Answer: B}

\end{frame}

% =========================================================
% Question 10
% =========================================================

\begin{frame}{Question 10}

The function $f$ is defined as

\[
f(x)=x^2+ax+b
\]

where $a$ and $b$ are positive integers.

If the minimum value of $f(x)$ is $7$, which is a possible value of $b$?

\begin{itemize}
    \item[A.] 6
    \item[B.] 7
    \item[C.] 8
    \item[D.] 9
\end{itemize}

\end{frame}

\begin{frame}{Answer to Question 10}

\textbf{Correct Answer: C}

\end{frame}

% =========================================================
% Question 11
% =========================================================

\begin{frame}{Question 11}

The average annual energy cost for a certain home is \$4,334.

The homeowner plans to spend \$25,000 to install a geothermal heating system.

The homeowner estimates that the average annual energy cost will then be \$2,712.

Which of the following inequalities can be solved to find $t$, the number of years after installation at which the total amount of energy cost savings will exceed the installation cost?

\begin{itemize}
    \item[A.] $25,000>(4,334-2,712)t$
    \item[B.] $25,000<(4,334-2,712)t$
    \item[C.] $25,000-4,334-2,712t$
    \item[D.]  $25000>\frac{4332}{2712}t$
\end{itemize}

\end{frame}

\begin{frame}{Answer to Question 11}

\textbf{Correct Answer: B}

\end{frame}

% =========================================================
% Question 12
% =========================================================

\begin{frame}{Question 12}

For the equation, $b$ and $c$ are constants.

For any real number $t$, the point

\[
\left(\frac{-2t+4}{3},t\right)
\]

lies on the graph of the equation

\[
bx+3y=c
\]

What is the value of $c$?

\begin{itemize}
    \item[A.] -2
    \item[B.] 4
    \item[C.] 6
    \item[D.] 12
\end{itemize}

\end{frame}

\begin{frame}{Answer to Question 12}

\textbf{Correct Answer: C}

\end{frame}

% =========================================================
% Question 13
% =========================================================

\begin{frame}{Question 13}

A circle in the $xy$-plane has equation

\[
(x-2.5)^2+(y-3.5)^2=225
\]

Line $p$ is tangent to the circle at point $(-6.5,C)$, where $C$ is a constant.

Which of the following could be the slope of line $p$?

\begin{itemize}
    \item[I.] $-\frac34$
    \item[II.] $\frac34$
    \item[III.] $\frac43$
\end{itemize}

\begin{itemize}
    \item[A.] I only
    \item[B.] I and II only
    \item[C.] I and III only
    \item[D.] I, II, and III
\end{itemize}

\end{frame}

\begin{frame}{Answer to Question 13}

\textbf{Correct Answer: B}

\end{frame}

% =========================================================
% Question 14
% =========================================================

\begin{frame}{Question 14}

The function $f$ is defined by

\[
f(x)=a^{x-b}
\]

where $a$ and $b$ are constants.

In the $xy$-plane, the graph of $y=f(x)$ passes through the points $(c,7)$ and $(2c,117)$, where $c$ is a constant.

What is the value of $b$?

\begin{itemize}
    \item[A.] 4
    \item[B.] 17
    \item[C.] 110
    \item[D.] 124
\end{itemize}

\end{frame}

\begin{frame}{Answer to Question 14}

\textbf{Correct Answer: A}

\end{frame}

% =========================================================
% Question 15
% =========================================================

\begin{frame}{Question 15}

The given equation, where $t$ is a positive constant, defines a circle in the $xy$-plane.

The radius of this circle is $\sqrt{41}$.

What is the value of $t$?

\[
2x^2+2y^2+tx+ty-78=0
\]

\begin{itemize}
    \item[A.] 1
    \item[B.] 4
    \item[C.] 7
    \item[D.] 8
\end{itemize}

\end{frame}

\begin{frame}{Answer to Question 15}

\textbf{Correct Answer: B}

\end{frame}

\end{document}